\section{Prototype dan Pewarisan (Tinjauan Mendalam)}

Setiap objek JavaScript memiliki \textbf{prototype}---objek lain yang menjadi ``template''-nya. Jika properti tidak ditemukan di objek, JavaScript mencari di prototype, lalu prototype-nya prototype, dan seterusnya (\textit{prototype chain}). Ini adalah mekanisme pewarisan di JavaScript; \code{class extends} hanyalah syntactic sugar di atasnya \cite{jsInfo}.

Properti \code{\_\_proto\_\_} (atau \code{Object.getPrototypeOf()}) mengakses prototype objek. Untuk constructor function, \code{F.prototype} adalah objek yang menjadi prototype dari instance yang dibuat dengan \code{new F()}. \code{Object.create(proto)} membuat objek baru dengan prototype tertentu. Memahami prototype chain membantu debug perilaku \code{instanceof}, method inheritance, dan \code{this} \cite{eloquentJs}.

\begin{konsep}
\textbf{Prototype Chain.} Ketika \code{obj.method()} dipanggil, JavaScript mencari \code{method} di: (1) objek itu sendiri, (2) prototype-nya, (3) prototype dari prototype, ... hingga \code{Object.prototype} (yang prototype-nya \code{null}). Ini menjelaskan mengapa semua objek memiliki \code{toString()}, \code{hasOwnProperty()}, dll.---mereka diwarisi dari \code{Object.prototype} \cite{mdnJsGuide}.
\end{konsep}

